\section{Numérique}

\begin{UPSTIexercice}{Conversion binaire vers décimal}
	\UPSTIetape Écrire un programme \texttt{binaire\_vers\_decimal.c} qui demande à l'utilisateur de saisir les bits d'un nombre binaire un par un, du poids fort vers le poids faible, jusqu'à ce qu'il entre \texttt{-1} pour arrêter la saisie. Le programme affichera la valeur décimale correspondante. Aucun tableau.
	Par exemple :
	\begin{lstlisting}[language=bash,style=console]
Entrez un bit (0 ou 1, -1 pour arreter) : 1
Entrez un bit (0 ou 1, -1 pour arreter) : 1
Entrez un bit (0 ou 1, -1 pour arreter) : 0
Entrez un bit (0 ou 1, -1 pour arreter) : 1
Entrez un bit (0 ou 1, -1 pour arreter) : -1
Valeur decimale : 13
\end{lstlisting}
\end{UPSTIexercice}

\begin{UPSTIcorrectionP}{Conversion binaire vers décimal}
	\begin{lstlisting}[language=c]
#include <stdio.h>

int main(void) {
    int bit;
    int valeur = 0;

    printf("Entrez un bit (0 ou 1, -1 pour arreter) : ");
    scanf("%d", &bit);

    while (bit != -1) {
        valeur = valeur * 2 + bit;
        printf("Entrez un bit (0 ou 1, -1 pour arreter) : ");
        scanf("%d", &bit);
    }

    printf("Valeur decimale : %d\n", valeur);

    return 0;
}
	\end{lstlisting}
\end{UPSTIcorrectionP}

\begin{UPSTIexercice}{Résolution d'un convertisseur analogique-numérique}
	\UPSTIetape Un convertisseur analogique-numérique (CAN) doit coder \texttt{N} valeurs distinctes. Écrire un programme \texttt{nombre\_bits.c} qui demande \texttt{N} et calcule, à l'aide d'une boucle, le plus petit nombre de bits \texttt{n} tel que $2^n \geq N$ (sans utiliser de fonction puissance).
	Par exemple :
	\begin{lstlisting}[language=bash,style=console]
Entrez le nombre de valeurs a coder : 1000
Nombre de bits necessaires : 10
\end{lstlisting}
\end{UPSTIexercice}

\begin{UPSTIcorrectionP}{Résolution d'un convertisseur analogique-numérique}
	\begin{lstlisting}[language=c]
#include <stdio.h>

int main(void) {
    int n;
    printf("Entrez le nombre de valeurs a coder : ");
    scanf("%d", &n);

    int bits = 0;
    long puissance = 1;
    while (puissance < n) {
        puissance *= 2;
        bits++;
    }

    printf("Nombre de bits necessaires : %d\n", bits);

    return 0;
}
	\end{lstlisting}
\end{UPSTIcorrectionP}

\begin{UPSTIexercice}{Bit de parité}
	\UPSTIetape Écrire un programme \texttt{parite.c} qui demande un entier positif \texttt{n} et compte, à l'aide d'une boucle, son nombre de bits à 1 (son \emph{poids de Hamming}), utile par exemple pour calculer un bit de parité en transmission numérique. Le programme affichera ce nombre et précisera si la parité est paire ou impaire. Aucun tableau.
	Par exemple :
	\begin{lstlisting}[language=bash,style=console]
Entrez un entier positif : 13
3 bit(s) a 1 (parite impaire)
\end{lstlisting}
\end{UPSTIexercice}

\begin{UPSTIcorrectionP}{Bit de parité}
	\begin{lstlisting}[language=c]
#include <stdio.h>

int main(void) {
    int n;
    printf("Entrez un entier positif : ");
    scanf("%d", &n);

    int reste = n;
    int nb_bits_a_1 = 0;
    while (reste != 0) {
        if (reste % 2 == 1) {
            nb_bits_a_1++;
        }
        reste /= 2;
    }

    printf("%d bit(s) a 1 ", nb_bits_a_1);
    if (nb_bits_a_1 % 2 == 0) {
        printf("(parite paire)\n");
    } else {
        printf("(parite impaire)\n");
    }

    return 0;
}
	\end{lstlisting}
\end{UPSTIcorrectionP}
